\section{Introduction}
\label{sec:intro}

World models that predict the future in the feature space of a frozen vision foundation model are cheap to train and plan with \citep{zhou2025dinowm,baldassarre2025dinoworld,assran2025vjepa2}. Their weakness is well known: a deterministic predictor trained with a squared loss returns the conditional mean, and when the future has several modes the mean lies between them, blurring the prediction and compounding over a rollout \citep{porcher2026flowwm,baldassarre2025dinoworld}. Stochastic feature-space world models fix this with diffusion or flow matching \citep{porcher2026flowwm,boduljak2025vfmf,walker2025generalist}, at the price of tens of network evaluations (NFE) per predicted step, which matters when a planner must imagine thousands of futures per decision. A growing line of work therefore builds one-step stochastic world models: drifting objectives \citep{lu2026driftworld}, mean-flow heads \citep{agibot2026geact}, terminal maps \citep{lee2026beckmann}, priming \citep{park2026priming}, or streaming flows with one velocity evaluation per step \citep{hsu2026vsf}.

The simplest one-step objective is older than all of these. Draw $m$ candidate futures from a noise-conditioned generator, find the one closest to the observed outcome, and pull only that one towards it. This is conditional implicit maximum likelihood estimation \citep[IMLE;][]{li2018imle,li2019diverse}, the ``min-of-$K$'' objective of explorative modeling \citep[XM;][]{gladstone2026xm}, and a sampled relative of winner-takes-all and multiple-choice learning \citep{guzman2012mcl,rupprecht2017mhp}. Every data point is explained by at least one sample, so mode collapse is discouraged by the objective itself, and inference is one forward pass. IMLE world models already outperform Gaussian ones on proprioceptive state \citep{aghabozorgi2026wimle}, and \citet{gladstone2026xm} conjecture that feature-space world models have few modes, so a small candidate budget should suffice. In the opposite direction, \citet{makabe2026esp} argue that a hard minimum over a finite number of candidates is not a strictly proper scoring rule and show that a one-step generator trained with the energy score covers a two-branch toy distribution better than the min-of-$K$ objective.

This paper asks a quantitative version of both claims: \emph{how many candidates does a one-step \nom{} world model need, when is no budget enough, and is a strictly proper score the better objective?} We answer with a conditional IMLE world model on frozen DINOv2/v3 patch tokens, trained and compared under matched parameters against deterministic regression, 50- and 2-step flow matching in the FlowWM recipe, two one-step flow generators (improved MeanFlow and Shortcut models), and an energy-score generator, on benchmarks whose futures can be enumerated so that mode coverage is measured rather than inferred.

One lens organizes the results. The \nom{} loss is, in expectation, the distortion incurred when the outcome is quantized with a \emph{random codebook} of $m$ independent draws from the generator. This correspondence is a sampled version of the known equivalence between winner-takes-all learning and conditional quantization \citep{rupprecht2017mhp,letzelter2024wta,perera2024amcl}, and its continuous limit is the Gaussian rate--distortion test channel \citep{shannon1959coding,cover2006elements}. We make no claim to new mathematics; the contribution is to show that this one correspondence predicts, quantitatively and before the experiments were run, where the candidate budget matters and where it cannot help.

\paragraph{Contributions.}
\begin{enumerate}
\item \textbf{A controlled test of the small-budget conjecture.} Under mode-like uncertainty (bouncing dynamics with $K$ enumerable outcomes, a Push-T variant with a hidden center of gravity), \nom{} with $m{=}16$ covers enumerated futures better than 50-step flow matching, the one-step flow generators and the energy-score generator at matched parameters (coverage $0.93$ vs.\ $0.79$ for FM-50 and $0.36$ for the energy score on non-repeating contexts; F1 error $7.2$ vs.\ $12.3$\,px on Bouncing Shapes), with about $30\times$ lower latency than FM-50 (Section~\ref{sec:results1}). On real robot video it does not, and Section~\ref{sec:results3} explains why.
\item \textbf{Outcome memorization and its incentive.} In contexts that never repeat, a conditional generator can reach zero training loss by memorizing the realized outcome of each training context, keyed on its appearance, and ignoring its noise input. We measure this with the share of samples that fall on the realized mode of training contexts (chance $1/K$), show that it is governed by the random codebook's miss \emph{distortion} $I(m)$ rather than by the miss probability $(1-1/K)^m$ (RMSE $0.031$ vs.\ $0.068$ over $K\in\{2,4,8,16\}$), and confirm seven of nine predictions recorded before the $K{=}2,16$ runs. The same accounting explains why the strictly proper energy score memorizes in this regime while \nom{} with $m{=}16$ does not: its incentive is half the expected outcome spread, independent of $m$ (Section~\ref{sec:results2}).
\item \textbf{Where no budget suffices.} The codebook view also predicts a flattening of mode probabilities with $m$, which we confirm within $0.01$ of the optimum derived by \citet{bennett2026freedom}, and an under-dispersion of $e_m/\sqrt{d}$ in $d$-dimensional continuous outcome spaces, which is the form of the Gaussian rate--distortion test channel. On BAIR robot pushing the samples spread along only one to three context-specific directions and no budget up to $m{=}64$ closes the gap to multi-step flow matching on distribution metrics (Section~\ref{sec:results3}).
\item \textbf{Rollouts and planning.} Autoregressive \nom{} rollouts stay as close to the token manifold as deterministic regression, while flow rollouts drift by a factor of four at $h{=}32$ regardless of the number of integration steps. Per-step matching on self-generated contexts, in the lineage of DaD and closest-among-top-$K$ rollouts \citep{venkatraman2015dad,zhang2025catk}, reduces long-horizon recall error by a third on stochastic Push-T; an ablation of the fed-back state separates its manifold benefit from its coverage benefit. Planning with a cross-entropy-method controller is $50\times$ faster than with FM-50 at equal success, but success does not improve even on a task where the model's branch-risk estimates are accurate (Section~\ref{sec:results4}).
\end{enumerate}

We also report what did not work: trajectory-level matching, multi-stage supervision in token space, variance-penalized planning and a counterfactual action-contrast term (Appendix~\ref{app:negative}). Benchmarks are synthetic or simulated except BAIR, models are small (8M and 21M parameters), and our flow-matching baselines are not at the scale of \citet{porcher2026flowwm}; Section~\ref{sec:discussion} states these limits.
